% Guangxi University Undergraduate Thesis/Design LaTeX Template
% Source: 广西大学资源环境与材料学院 论文审查心得及注意事项 + 本科毕业论文参考模板
% Compile with XeLaTeX.

\documentclass[UTF8,zihao=-4,linespread=1.0,twoside,scheme=chinese]{ctexrep}

% ===== Page Geometry =====
% Spec: A4, top 2.54cm, bottom 2.54cm, left 2.2cm, right 2.2cm
% Header 2cm from top, footer 1.75cm from bottom
\usepackage[a4paper,
  left=2.2cm, right=2.2cm,
  top=2.54cm, bottom=2.54cm,
  headheight=0.8cm, headsep=1.2cm,
  footskip=0.79cm]{geometry}

% ===== Fonts =====
% Spec: 宋体(SimSun) for body, 黑体(SimHei) for headings, Times New Roman for English
% macOS: Songti SC / Heiti SC / Kaiti SC; Windows: SimSun / SimHei / KaiTi
\setmainfont{Times New Roman}
\setCJKmainfont{Songti SC}[
  BoldFont=Heiti SC,
  ItalicFont=Kaiti SC,
  BoldItalicFont=Heiti SC]
\setCJKsansfont{Heiti SC}[BoldFont=Heiti SC]
\setCJKmonofont{Kaiti SC}[BoldFont=Heiti SC]
\setmonofont{Courier New}

% ===== Line Spacing =====
% Spec: fixed 20pt throughout body text
% 小四(12pt) with 20pt baselineskip
\usepackage{setspace}
\renewcommand{\normalsize}{\fontsize{12pt}{20pt}\selectfont}
\setstretch{1.0}

% ===== First-line Indent =====
% Spec: 2 Chinese characters
\usepackage{indentfirst}
\setlength{\parindent}{2em}

% ===== Headers and Footers =====
\usepackage{fancyhdr}
\fancyhf{}
% Header: left "广西大学本科毕业论文", right paper title, 隶书小四(12pt)
% 隶书 not available on macOS; use Kaiti SC as substitute
\fancyhead[LE]{\kaiti\zihao{-4} 广西大学本科毕业论文}
\fancyhead[RO]{\kaiti\zihao{-4} $if(title)$$title$$endif$}
\fancyhead[RE]{\kaiti\zihao{-4} 广西大学本科毕业论文}
\fancyhead[LO]{\kaiti\zihao{-4} $if(title)$$title$$endif$}
% Footer: page number centered, Times New Roman 五号(10.5pt)
\fancyfoot[C]{\zihao{5}\thepage}
\renewcommand{\headrulewidth}{0.4pt}
\renewcommand{\footrulewidth}{0pt}

\fancypagestyle{plain}{
  \fancyhf{}
  \fancyfoot[C]{\zihao{5}\thepage}
  \renewcommand{\headrulewidth}{0pt}
}

% ===== Heading Formatting =====
% Spec: Level 1 小二号黑体(18pt) centered, blank line before/after
% Level 2 小三号黑体(15pt) left-aligned, no extra space
% Level 3 四号黑体(14pt) left-aligned, no extra space
% Level 4 小四号黑体(12pt) left-aligned, no extra space
\usepackage{titlesec}
\titleformat{\chapter}[hang]
  {\centering\heiti\zihao{-2}\bfseries}
  {第\thechapter 章}{1em}{}
\titlespacing{\chapter}{0pt}{20pt}{20pt}

\titleformat{\section}
  {\heiti\zihao{-3}\bfseries\raggedright}
  {\thesection}{1em}{}
\titlespacing{\section}{0pt}{0pt}{0pt}

\titleformat{\subsection}
  {\heiti\zihao{4}\bfseries\raggedright}
  {\thesubsection}{1em}{}
\titlespacing{\subsection}{0pt}{0pt}{0pt}

\titleformat{\subsubsection}
  {\heiti\zihao{-4}\bfseries\raggedright}
  {\thesubsubsection}{1em}{}
\titlespacing{\subsubsection}{0pt}{0pt}{0pt}

% ===== Table of Contents =====
% Spec: "目 录" 三号黑体(16pt) centered, content 小四(12pt), 1.25x line spacing
\usepackage{tocloft}
\renewcommand{\cfttoctitlefont}{\hfill\heiti\zihao{3}}
\renewcommand{\cftaftertoctitle}{\hfill}
\renewcommand{\contentsname}{目\ \ 录}
\renewcommand{\cftchapfont}{\zihao{-4}}
\renewcommand{\cftsecfont}{\zihao{-4}}
\renewcommand{\cftsubsecfont}{\zihao{-4}}
\renewcommand{\cftchappagefont}{\zihao{-4}}
\renewcommand{\cftsecpagefont}{\zihao{-4}}
\renewcommand{\cftsubsecpagefont}{\zihao{-4}}
% TOC line spacing: 1.25x
\renewcommand{\cftbeforechapskip}{0pt}
\renewcommand{\cftbeforesecskip}{0pt}
\renewcommand{\cftbeforesubsecskip}{0pt}
% Secondary heading indent 2 chars, tertiary 4 chars
\renewcommand{\cftsecindent}{2em}
\renewcommand{\cftsubsecindent}{4em}

% ===== Captions =====
% Spec: 五号(10.5pt) bold, figure below/table above, centered, no trailing punctuation
\usepackage{caption}
\usepackage[font=bf,labelsep=quad]{caption}
\captionsetup{
  font={bf,small},
  labelfont={bf},
  labelsep=quad,
  justification=centering,
  singlelinecheck=false}

% ===== Three-line Tables =====
% Spec: top/bottom 1.5pt, header separator 0.5pt
\usepackage{booktabs}
\renewcommand{\heavyrulewidth}{1.5pt}
\renewcommand{\lightrulewidth}{0.5pt}

% ===== Figures and Graphics =====
\usepackage{graphicx}
\graphicspath{{figures/}{images/}{./}}

% ===== Tables =====
\usepackage{longtable}
\usepackage{array}

% ===== Mathematics =====
\usepackage{amsmath}

% ===== Hyperlinks =====
\usepackage{hyperref}
\usepackage{xurl}
\hypersetup{
  colorlinks=true,
  linkcolor=black,
  citecolor=black,
  urlcolor=black,
  breaklinks=true}

% ===== Labels =====
\renewcommand{\bibname}{参考文献}
\renewcommand{\figurename}{图}
\renewcommand{\tablename}{表}

% ===== Pandoc Compatibility =====
\providecommand{\tightlist}{\setlength{\itemsep}{0pt}\setlength{\parskip}{0pt}}
\providecommand{\pandocbounded}[1]{#1}
% Pandoc longtable uses \LTcaptype{none} for unnumbered tables
\expandafter\ifx\csname c@none\endcsname\relax
  \newcounter{none}
\fi

$if(csl-refs)$
\newlength{\cslhangindent}
\setlength{\cslhangindent}{1.5em}
\newenvironment{CSLReferences}[2]{\begin{thebibliography}{99}}{\end{thebibliography}}
\newcommand{\CSLBlock}[1]{#1\par}
\newcommand{\CSLLeftMargin}[1]{\makebox[\cslhangindent][l]{#1}}
\newcommand{\CSLRightInline}[1]{#1}
\newcommand{\CSLIndent}[1]{\hspace{\cslhangindent}#1}
$endif$

% ===== Commands =====
\newcommand{\kaiti}{\CJKfamily{zhkai}}

\begin{document}

% ===== Title Page =====
% Spec: 黑体一号(26pt) centered, other info 宋体四号(14pt) centered
\begin{titlepage}
\centering
\vspace*{2.0cm}
{\heiti\zihao{1} 本科毕业论文 \par}
\vspace{0.8cm}
{\heiti\zihao{1} $if(title)$$title$$else$课题名称$endif$ \par}
\vfill
{\songti\zihao{4}
\begin{tabular}{rl}
学院： & $if(college)$$college$$else$资源环境与材料学院$endif$ \\
专业： & $if(major)$$major$$else$专业名称$endif$ \\
班级： & $if(class-name)$$class-name$$else$班级$endif$ \\
学号： & $if(student-id)$$student-id$$else$学号$endif$ \\
姓名： & $if(author)$$author$$else$姓名$endif$ \\
指导老师： & $if(advisor)$$advisor$$else$指导教师$endif$ \\
\end{tabular}
\par}
\vspace{1.5cm}
{\songti\zihao{4} $if(date)$$date$$else$二〇二六年五月$endif$ \par}
\end{titlepage}

% ===== Front Matter: Single-sided, Roman numerals =====
\pagenumbering{Roman}
\setcounter{page}{1}

% ===== Chinese Abstract =====
% Spec: Chinese title 黑体三号(16pt) centered, "摘 要" 黑体三号(16pt) centered
% Body 宋体小四(12pt) fixed 20pt, 2-char indent, 400-600 chars
% Keywords: 黑体小四, content 宋体小四
$if(abstract-zh)$
\cleardoublepage
\begin{center}
{\heiti\zihao{3} $if(title)$$title$$endif$ \par}
\vspace{20pt}
{\heiti\zihao{3} 摘\ \ 要 \par}
\end{center}
\vspace{20pt}
$abstract-zh$

\vspace{20pt}
\noindent{\heiti\zihao{-4} 关键词：}{\songti\zihao{-4} $if(keywords-zh)$$keywords-zh$$endif$}
$endif$

% ===== English Abstract =====
% Spec: Title ALL CAPS TNR 三号(16pt) bold centered
% "ABSTRACT" 三号(16pt) bold centered
% Body TNR 小四(12pt) fixed 20pt, 4-char indent per paragraph
% "KEYWORDS" bold, left-aligned, no indent
$if(abstract-en)$
\cleardoublepage
\begin{center}
{\bfseries\zihao{3} \MakeUppercase{$if(title-en)$$title-en$$else$$title$$endif$} \par}
\vspace{20pt}
{\bfseries\zihao{3} ABSTRACT \par}
\end{center}
\vspace{20pt}
$abstract-en$

\vspace{20pt}
\noindent{\bfseries KEYWORDS: } $if(keywords-en)$$keywords-en$$endif$
$endif$

% ===== Table of Contents =====
% Spec: "目 录" 三号黑体(16pt) centered, content 小四(12pt), 1.25x spacing
% Secondary headings indent 2 chars, tertiary 4 chars
\cleardoublepage
{
\renewcommand{\baselinestretch}{1.25}
\tableofcontents
}

% ===== Main Body: Double-sided, Arabic numerals =====
\cleardoublepage
\pagenumbering{arabic}
\setcounter{page}{1}
\pagestyle{fancy}

$body$

$if(nocite-ids)$
\nocite{$for(nocite-ids)$$it$$sep$, $endfor$}
$endif$

% ===== References =====
$if(bibliography)$
\clearpage
$bibliography$
$endif$

% ===== Acknowledgements =====
% Spec: 300-600 chars, sincere and concise
$if(acknowledgements)$
\clearpage
\chapter*{致谢}
\addcontentsline{toc}{chapter}{致谢}
$acknowledgements$
$endif$

\end{document}
